% Drawing 154/03-15 -- folder 154, batch 03, archive page 55 (PDF page 56).
% Transcription: figures au crayon et à l'encre : un disque au crayon, à la périphérie duquel sont accolés un pentagone et un hexagone orientés (petites flèches courbes), un troisième polygone esquissé au crayon à leur gauche, deux points $0$ et $1$ sur le bord ; à droite, un segment vertical portant les points $s_1$, $s_0$, $s_{n-1}$
% Kind: disk with oriented polygons on its rim; segment with marked points. Method 2 (schematic redraw from the facsimile). Status: ok.
% Notes: Pencil circle; on its upper rim an inked pentagon and an inked hexagon side by side, sharing the hexagon's left side, a curved arrow inside each (clockwise in the pentagon, beside a figure 1; counter-clockwise in the hexagon, round a figure 0), a third polygon sketched in pencil on their left (an open broken line and a separate short stroke); a full point at the foot of the pentagon, an open point and the figures 0 and 1 under the hexagon; an inked stretch of the rim joining the feet of the two polygons; a pencil arrowhead on the rim at the upper left; four faint pencil chords inside the circle (two long, two short), a short pencil stroke by the figure 0 and a short dash outside the rim at the right. The arrowheads on three sides of the polygons (the pentagon's left side, the hexagon's left side and its lower right side, all pointing to the rim) are drawn. The lines of formulas written across the circle belong to the text. Right: the vertical segment with an open point at its top, two ticks, a small pencil hook under the second, the open point s_1, a tick, a thick stretch ending in a tick, the full point s_0 inside a pencil oval (its left arc and two strokes on its right), a tick, a second thick stretch between two short cross strokes, the full point s_{n-1}; the pencil strokes beside it (a horizontal stroke left of s_0, a vertical stroke left of the lower thick stretch, a horizontal stroke right of s_{n-1}, a bracket on the right below it). Redrawn 2026-11-12 (issue 128). 2026-11-14, Claude Opus 5.5, second pass (issue 147): redrawn stroke by stroke from the 400 dpi colour crop, coordinates read off a unit grid; the circle is drawn as a smooth curve through 27 points read on its ink, the polygons, the curved arrows, the points and labels moved onto their ink (the points and labels were 2 units low), the segment on the right redrawn with its ticks at their places (two of them were 15 units off). Overlay: every stroke lies on its ink (mean distance under 0.5 unit; only the two faintest pencil strokes reach 1.5 to 2 units at their ends); the small hooks at the ends of the cross strokes of the lower thick stretch are drawn as plain ticks. Not compiled (no TeX in the session). Not checked by a reader.
% Coordinates: 1 unit = 1 mm = 1/100 of the crop width (bbox in page fractions: [0.22, 0.085, 0.42, 0.275]).
% Redrawn by Claude (claude-fable-5-2), issue #128, 2026-11-12; not checked by a human; not compiled (no TeX in the session).
% Redrawn stroke by stroke by Claude (claude-opus-5-5) 2026-11-14 (issue #147, second pass): overlaid on the 400 dpi colour facsimile on a unit grid; not compiled (no TeX in the session). Not checked by a human. Status: ok.
% Rights: reproduces part of the work -- behind RENDER_TRANSCRIPTS, never published by itself.
\documentclass[tikz,border=3pt]{standalone}
\usepackage{amsmath,amssymb}
\usetikzlibrary{patterns,arrows.meta}
\tikzset{pencil/.style={black!45}, struck/.style={black!35}, mpt/.style={circle,draw,inner sep=1.3pt}, pt/.style={circle,fill,inner sep=0.9pt}, opt/.style={circle,draw,fill=white,inner sep=1.1pt}}
\begin{document}
\begin{tikzpicture}[x=1mm,y=1mm,line cap=round,line join=round,font=\scriptsize]
\draw[pencil] plot[smooth cycle, tension=0.6] coordinates {(52.3,-69.9) (57.5,-73.8) (61.8,-78.4) (64.4,-83) (66.3,-90) (66.6,-98) (65.3,-106) (61,-114.4) (55,-120.7) (48,-124.6) (40,-126.7) (34,-127.2) (26,-126.6) (18,-123.6) (12,-118.6) (7.5,-111.5) (5.5,-104) (5.4,-98) (6.7,-90) (9.5,-82.6) (14,-76.8) (19.5,-73.2) (24,-70.6) (28.3,-69.1) (32.3,-68.1) (41.7,-66.4) (47,-67.2)};
\draw[pencil] (19.5,-73.2) -- (23,-73.9);
\draw (28.3,-69.1) -- (32.3,-68.1) -- (36.5,-67.2) -- (41.7,-66.4) -- (47,-67.2) -- (52.3,-69.9);
\draw (39,-56.5) -- (29.6,-50.3) -- (25,-57.6) -- (28.3,-69.1);
\draw (41.7,-66.4) -- (41.3,-57) -- (47.3,-50) -- (52,-50.9) -- (57.3,-51.3) -- (59.1,-60.2) -- (52.3,-69.9);
\draw (26.1,-64.3) -- (28.3,-69.1) -- (28.1,-63.8);
\draw (40.2,-63.4) -- (41.7,-66.4) -- (43.8,-61.6);
\draw (53.4,-65.9) -- (52.3,-69.9) -- (56.4,-67);
\draw plot[smooth, tension=0.6] coordinates {(31.2,-55.5) (34.6,-57) (35.9,-60.3) (34.4,-63) (31.6,-62.4)};
\draw (32.6,-61.2) -- (31.6,-62.4) -- (32.6,-64.1);
\draw plot[smooth, tension=0.6] coordinates {(46.4,-60.8) (47.6,-63.4) (49.8,-64.2) (52.2,-62.5) (52.8,-59.4) (51.9,-56.8) (50,-55.5)};
\draw (50.2,-58.6) -- (50,-55.5) -- (52.4,-56.7);
\node[inner sep=1pt] at (30.4,-58.3) {\small $1$};
\node[inner sep=1pt] at (48.4,-59.8) {\small $0$};
\draw[pencil] (15.2,-55.7) -- (10.5,-62.7) -- (16.5,-74.1);
\draw[pencil] (16,-55.6) -- (22.7,-58.6);
\draw[pencil] (23.8,-74.5) -- (20.2,-104);
\draw[pencil] (42.3,-77.5) -- (39.4,-107);
\draw[pencil] (31.4,-82) -- (28.8,-100);
\draw[pencil] (54.1,-92) -- (53.4,-99.7);
\node[pt] at (32.3,-68.1) {};
\draw[pencil] (45.5,-68.1) -- (46.6,-70.6);
\draw[pencil] (66,-74) -- (69,-73.3);
\node[opt] at (40.9,-71.3) {};
\node[inner sep=1pt] at (47.4,-73.4) {\small $0$};
\node[inner sep=1pt] at (51.6,-73.1) {\small $1$};
\draw (89.4,-8.3) -- (90.4,-128);
\node[opt] at (89.4,-8.3) {};
\draw (88.2,-15.6) -- (91.5,-15.6);
\draw (88,-32) -- (91.3,-32);
\draw[pencil] (88.1,-32.3) -- (88.9,-34.1) -- (89.6,-34);
\node[opt] at (90,-38.2) {};
\draw (87.7,-49.3) -- (91.6,-49.3);
\draw[line width=1.2pt] (90,-49.3) -- (90.1,-62.5);
\draw (87.7,-62.5) -- (90.8,-62.5);
\draw[pencil] plot[smooth, tension=0.6] coordinates {(87.4,-49.7) (86,-55) (85.6,-62) (86,-69) (87,-74.4)};
\draw[pencil] (93.2,-51) -- (93.2,-60.5);
\draw[pencil] (92.9,-66) -- (91.6,-74);
\node[pt] at (89.6,-70) {};
\draw[pencil] (80.3,-70.6) -- (87.6,-70.6);
\draw (88.2,-75.5) -- (91.2,-75.5);
\draw[line width=1.2pt] (90,-77) -- (90.1,-87);
\draw (88.6,-78.5) -- (92.5,-78.5);
\draw (88.6,-87.5) -- (92.5,-87.5);
\node[pt] at (88.6,-92) {};
\draw[pencil] (89,-92.3) -- (95.5,-92.3);
\draw[pencil] (79.5,-76.5) -- (79.5,-89.5);
\draw[pencil] (91.5,-88) -- (91.4,-100.5) -- (90.4,-101);
\node[inner sep=1pt] at (97.5,-40) {\small $s_1$};
\node[inner sep=1pt] at (95.5,-70.5) {\small $s_0$};
\node[inner sep=1pt] at (82,-94.5) {\small $s_{n-1}$};
\end{tikzpicture}
\end{document}
